\documentclass[twoside]{article}
\usepackage[preprint]{aistats2027}
\usepackage[round,authoryear]{natbib}
\usepackage{amsmath,amssymb,amsthm}
\usepackage{mathtools}
\usepackage{microtype}
\usepackage{booktabs}
\usepackage{graphicx}
\usepackage{url}
\renewcommand{\bibname}{References}
\renewcommand{\bibsection}{\subsubsection*{\bibname}}
\bibliographystyle{apalike}
\newtheorem{theorem}{Theorem}
\newtheorem{lemma}[theorem]{Lemma}
\newtheorem{proposition}[theorem]{Proposition}
\newtheorem{corollary}[theorem]{Corollary}
\newcommand{\E}{\mathbb E}
\newcommand{\R}{\mathbb R}
\newcommand{\D}{\mathsf D}
\newcommand{\TV}{d_{\mathrm{TV}}}
\newcommand{\Reg}{\operatorname{Reg}}
\newcommand{\supp}{\operatorname{supp}}
\begin{document}
\runningauthor{}
\twocolumn[
\aistatstitle{Optimal-Dimension U-Calibration by Bayesian Bootstrap}
\aistatsauthor{Pahan Dewasurendra}
\aistatsaddress{Johns Hopkins University}
]
\begin{abstract}
U-calibration asks one online probability forecaster to have low regret for every bounded proper loss, including losses unknown when the forecasts are made. For nontrivial $K\geq2$, the optimal worst-case rate is $\Theta(\!\sqrt{\min\{K,T\}T}\,)$, while the smooth-loss class has minimax time rate $\Theta(\log T)$. A recent simultaneous algorithm attains both time rates but leaves a polynomial gap in $K$. We show that the classical Bayesian bootstrap closes this gap. At round $t$, independently assign exponential weights to the past outcomes and predict their normalized weighted histogram. Equivalently, if $c_{t-1}$ is the vector of past outcome counts, sample $P_t\sim\operatorname{Dirichlet}(c_{t-1})$ on the observed face of the simplex. For every outcome sequence with terminal counts $n_1,\ldots,n_K$, observed support size $m$, and every proper loss in $[-1,1]$, the expected regret is at most \[ 6\sum_{i=1}^K\sqrt{n_i}\leq6\sqrt{mT}\leq6\sqrt{\min\{K,T\}T}. \] For every $\beta$-smooth loss, the same forecasts have regret at most the regret of empirical follow-the-leader plus $(\beta/2)\log T$, and therefore $O(\beta\log T)$ for bounded smooth proper losses. The proof uses three elementary facts: a one-count Dirichlet update has total variation $O(1/\sqrt{c_i})$, size-biasing an exponential weight adds an independent exponential weight, and weighted be-the-leader holds simultaneously for all proper losses. A class-aggregation reduction gives the matching dimension-capped lower bound. The algorithm is loss-agnostic, horizon-free, proper, and minimax-optimal for the worst bounded loss while retaining the minimax-optimal smooth-class time rate.
\end{abstract}

\section{Introduction}

A probability forecast is often consumed by a decision maker whose utility is unknown to the forecaster. Proper losses formalize this setting. If outcomes follow a distribution $p$, reporting $p$ minimizes expected loss. U-calibration asks for one sequence of online forecasts that has low regret under every bounded proper loss \citep{kleinberg2023}. This gives a loss-agnostic guarantee to all downstream agents represented by those losses.

For $K$ outcomes and horizon $T\geq12K$, the minimax worst-loss rate is $\Theta(\sqrt{KT})$ \citep{luo2024}. The upper bound follows by adding independent geometric numbers of hallucinated copies of each outcome before taking their empirical frequency. This perturbation is optimal for the worst bounded proper loss, but it is too large for easier losses. In particular, it incurs $\Omega(\sqrt T)$ regret even for squared loss, whose optimal regret is $\Theta(\log T)$ \citep{frongillo2026}. We also show that the all-regime worst-loss rate is $\Theta(\!\sqrt{\min\{K,T\}T}\,)$ by lifting the same lower bound through outcome aggregation.

\citet{frongillo2026} recently showed that this conflict is not inherent. Their self-concordant perturbation remains inside the probability simplex and simultaneously obtains logarithmic regret for smooth proper losses. For general bounded proper losses, however, its regret is $\widetilde O(K^{5/4}\sqrt T)$ rather than the optimal $O(\sqrt{KT})$. They explicitly leave optimal dependence on $K$ open. A Gaussian construction has the desired dimension dependence but can predict outside the simplex, where a general proper loss is not defined.

We close this dimension gap with a classical randomization from statistics. The Bayesian bootstrap draws independent exponential weights on the observations and recomputes the statistic of interest \citep{rubin1981}. Applied to a categorical sample, its weighted empirical distribution is Dirichlet with parameters equal to the observed counts. Using a fresh bootstrap replicate at every round gives our forecaster.

The result is unexpectedly direct. If the current label has appeared $j$ times, adding its latest occurrence changes the Dirichlet law by total variation $O(j^{-1/2})$. Summing over occurrences gives $O(\sum_i\sqrt{n_i})$. This controls the usual stability term. The less obvious be-the-perturbed-leader term is handled by exponential size bias. Multiplying by the current exponential weight changes that weight from $\operatorname{Gamma}(1,1)$ to $\operatorname{Gamma}(2,1)$, which is the same as adding one independent exponential copy. Weighted be-the-leader then compares a one-count-shifted Dirichlet law to the best fixed forecast. A second application of the same total-variation bound completes the proof.

Smooth adaptation comes from another basic property of the Bayesian bootstrap. A Dirichlet draw centered at the empirical frequency is unbiased, with total squared variance at most $1/t$. Smoothness therefore charges at most $\beta/(2t)$ relative to empirical follow-the-leader. The harmonic sum preserves logarithmic regret.

\paragraph{Contributions.} Our main theorem gives a single horizon-free algorithm with
\[ \Reg_\ell\leq6\sum_i\sqrt{n_i}\leq6\sqrt{mT}\leq6\sqrt{\min\{K,T\}T} \] for every bounded proper loss $\ell$. The first inequality is support and frequency adaptive, and a class-aggregation corollary proves that the final rate is minimax-optimal in every dimension regime. For every bounded $\beta$-smooth proper loss, the same algorithm has $O(\beta\log T)$ regret, which matches the smooth-class time lower bound. Unlike prior optimal-worst-case algorithms, the perturbation scale is not tuned to $T$. Unlike prediction-space Gaussian noise, every forecast belongs to the simplex.

\paragraph{Scope.} We follow the simultaneous expected-regret convention of \citet{frongillo2026}: one loss-agnostic algorithm satisfies a bound for every fixed loss. This is also called pseudo U-calibration, $\sup_\ell\E[\Reg_\ell]$, in \citet{luo2024}. We do not claim a bound on $\E[\sup_\ell\Reg_\ell]$. The self-concordant method also gives logarithmic regret for some nonsmooth losses that are smooth relative to the log barrier. Our theorem does not subsume that extension.

\section{Setting and Prior Bounds}

Let $[K]=\{1,\ldots,K\}$ and let \[ \Delta_K=\left\{p\in\R_+^K:\sum_{i=1}^Kp_i=1\right\}. \] At round $t$, the forecaster outputs $P_t\in\Delta_K$ and observes $y_t\in[K]$. Write $e_i$ for the $i$th standard basis vector, $c_t=\sum_{s=1}^t e_{y_s}$ for the count vector, and $n_i=c_{T,i}$. Let $m=|\{i:n_i>0\}|$ be the number of observed outcomes.

A loss $\ell:\Delta_K\times[K]\to[-1,1]$ is \emph{proper} if, for every $q\in\Delta_K$, \[ q\in\arg\min_{p\in\Delta_K}\sum_{i=1}^Kq_i\ell(p,i). \] For a fixed outcome sequence, expected regret is \[ \Reg_\ell(T)=\E\!\left[\sum_{t=1}^T\ell(P_t,y_t)\right]-\min_{p\in\Delta_K}\sum_{t=1}^T\ell(p,y_t). \] Properness implies that the comparator is the terminal empirical frequency $c_T/T$. Expectations are only over the forecaster.

The main statement is for an oblivious sequence. It also holds when $y_t$ is selected from the history before the fresh random draw at time $t$. This is the standard nonanticipating adaptive protocol. Section~S3 of the supplement gives the conditional argument and relates it to the usual fresh-perturbation reduction \citep{hutter2005}.

A differentiable loss is $\beta$-smooth if for all $p,q\in\Delta_K$ and $y\in[K]$, \[ \ell(p,y)\leq \ell(q,y)+\langle\nabla\ell(q,y),p-q\rangle+\frac\beta2\|p-q\|_2^2. \] We take $\beta\geq1$, which loses nothing because a loss with a smaller constant is also one-smooth. Empirical follow-the-leader predicts $Q_1=P_1$ and $Q_t=c_{t-1}/(t-1)$ for $t\geq2$. Prior work gives $\Reg_\ell^{\mathrm{FTL}}(T)=O(\beta\log T)$ for bounded $\beta$-smooth proper losses \citep{luo2024,frongillo2026}.

Table~\ref{tab:comparison} compares simultaneous methods. Geometric count perturbations obtain the optimal worst-loss rate but not smooth adaptation. Self-concordant noise adapts, but has a larger $K$ exponent in the worst case. The Bayesian bootstrap attains both target rates. Concurrent calibeating work reduces calibeating to regret minimization \citep{chen2026} or gives loss-based regularization for Lipschitz and Tsallis families \citep{fichtl2026}. Those results do not provide a worst-case-optimal learner for all bounded proper losses that also adapts to smooth losses.

\begin{table*}[t]
\caption{Per-loss expected regret for one loss-agnostic multiclass forecaster. Logarithmic entries concern bounded $\beta$-smooth proper losses. The count-adaptive bound implies the stated worst-case rate.}
\label{tab:comparison}
\centering
\small
\resizebox{0.98\textwidth}{!}{\begin{tabular}{lccc}
\toprule
Method & Every bounded proper loss & Every smooth proper loss & Horizon-free and proper \\
\midrule
Geometric count FTPL \citep{luo2024} & $O(\sqrt{KT})$ & $\Omega(\sqrt T)$ can occur & No, yes \\
Self-concordant noise \citep{frongillo2026} & $\widetilde O(K^{5/4}\sqrt T)$ & $O(\beta\log T+\beta\sqrt K\log K)$ & No, yes \\
Gaussian prediction noise \citep{frongillo2026} & $O(\sqrt{KT\log T})$ & $O(\beta\log T)$ & No, no \\
Bayesian bootstrap, this paper & $6\sum_i\sqrt{n_i}\leq6\sqrt{\min\{K,T\}T}$ & $O(\beta\log T)$ & Yes, yes \\
\bottomrule
\end{tabular}}
\end{table*}

\paragraph{Adjacent benchmarks.} Several recent lines ask stronger questions under different comparators. Proper-calibeating controls calibration and refinement uniformly over bounded proper losses relative to reference forecasts \citep{fosterhart2026}. Recalibration balances calibration with excess loss relative to binary hint sequences \citep{hu2026recalibration}. Omniprediction and decision-swap regret use feature-dependent comparator classes or loss-dependent postprocessing \citep{okoroafor2025,lu2025,hu2026blackwell}. Most recently, swap-agnostic learning lets a binary comparator hypothesis vary with the learner's prediction bucket \citep{okoroafor2026}. None of these results supplies the optimal finite-horizon multiclass regret against the best fixed distribution for every bounded proper loss. Conversely, our best-fixed-distribution guarantee does not imply calibration, hinted excess-loss control, or prediction-conditioned comparator regret.

\section{Bayesian-Bootstrap Forecaster}

For a nonzero count vector $c\in\mathbb N_0^K$, let $\D_c$ denote the Dirichlet law on its observed face. Formally, draw independent $G_i\sim\operatorname{Gamma}(c_i,1)$ for $c_i>0$, set $G_i=0$ when $c_i=0$, and return $G/\sum_iG_i$. This definition includes lower-dimensional faces. If only one coordinate is positive, $\D_c$ is a point mass.

\paragraph{Algorithm.} Choose any fixed $P_1\in\Delta_K$. At each round $t\geq2$, independently sample
\[ P_t\sim\D_{c_{t-1}}. \tag{BB} \] After observing $y_t$, update $c_t=c_{t-1}+e_{y_t}$. Equivalently, draw independent $W_{t,s}\sim\operatorname{Exp}(1)$ for $s<t$ and predict \[ P_t=\frac{\sum_{s<t}W_{t,s}e_{y_s}}{\sum_{s<t}W_{t,s}}. \] The equivalence follows because sums of exponential variables are gamma and normalized independent gammas are Dirichlet. It is also exactly the Bayesian bootstrap after grouping repeated categorical observations \citep{rubin1981}.

The count representation samples at most $m$ gamma variables per round and stores $K$ integer counts. It does not use $T$, $\ell$, or $\beta$.

\begin{theorem}[Worst-loss guarantee]
\label{thm:worst}
For every $K,T$, every outcome sequence, and every proper $\ell:\Delta_K\times[K]\to[-1,1]$, the Bayesian-bootstrap forecaster satisfies \[ \Reg_\ell(T)\leq6\sum_{i:n_i>0}\sqrt{n_i}\leq6\sqrt{mT}\leq6\sqrt{\min\{K,T\}T}. \] The same inequalities hold against a nonanticipating adaptive adversary in expectation over the realized outcome path.
\end{theorem}

\begin{theorem}[Smooth transfer]
\label{thm:smooth}
Let $\ell$ be a $\beta$-smooth proper loss and let $\Reg_\ell^{\mathrm{FTL}}(T)$ be the pathwise regret of empirical follow-the-leader initialized at the same $P_1$. Then \[ \Reg_\ell(T)\leq \Reg_\ell^{\mathrm{FTL}}(T)+\frac\beta2\sum_{t=2}^T\frac1t. \] Consequently, every bounded $\beta$-smooth proper loss has $\Reg_\ell(T)=O(\beta\log T)$ under the same forecasts used in Theorem~\ref{thm:worst}.
\end{theorem}

\begin{corollary}[All-regime minimax rate]
\label{cor:minimax}
Let $\mathcal L_K$ be the proper losses from $\Delta_K\times[K]$ to $[-1,1]$. For $K\geq2$ and $T\geq24$, \[ \inf_{\mathcal A}\sup_{\substack{y_{1:T}\in[K]^T\\ \ell\in\mathcal L_K}} \Reg^{\mathcal A}_\ell(T) =\Theta\!\left(\sqrt{\min\{K,T\}T}\right), \] where the infimum is over randomized forecasting algorithms. The upper bound is horizon-free. Section~S5 proves the lower bound by aggregating $K$ outcomes onto $r=\min\{K,\lfloor T/12\rfloor\}$ outcomes and invoking the lower bound of \citet{luo2024}.
\end{corollary}

For a realized randomized forecast sequence, write \[ \widehat\Reg_\ell(T)=\sum_{t=1}^T\ell(P_t,y_t)-\min_{p\in\Delta_K}\sum_{t=1}^T\ell(p,y_t). \] The fresh draws also give a simultaneous guarantee over a finite collection of losses.

\begin{corollary}[Finite loss families]
\label{cor:finite}
Let $\mathcal L_0$ be a finite family of $M$ proper losses in $[-1,1]$. Against an oblivious or nonanticipating adaptive adversary, with probability at least $1-\delta$, \[ \max_{\ell\in\mathcal L_0}\widehat\Reg_\ell(T) \leq6\sum_i\sqrt{n_i}+\sqrt{2T\log(M/\delta)}. \tag{1} \] For a fixed oblivious sequence, \[ \E\max_{\ell\in\mathcal L_0}\widehat\Reg_\ell(T) \leq6\sum_i\sqrt{n_i}+\sqrt{2T\log M}. \tag{2} \] Thus the stronger expected-supremum notion is controlled for every finite loss family.
\end{corollary}

Corollary~\ref{cor:minimax} makes the bounded-loss guarantee minimax-optimal in every dimension regime. The $\Omega(\log T)$ squared-loss lower bound of \citet{luo2024} makes the smooth-class time rate optimal. The first inequality of Theorem~\ref{thm:worst} is stronger when only a few classes occur or the terminal histogram is unbalanced.

\section{One-Count Dirichlet Stability}

The proof rests on an elementary update lemma. We use $\TV(\mu,\nu)=\tfrac12\int|d\mu-d\nu|$.

\begin{lemma}[One-count update]
\label{lem:tv}
Let $c\neq0$, $N=\sum_kc_k$, and $c_i>0$. Then \[ \TV(\D_c,\D_{c+e_i}) \leq\frac12\sqrt{\frac{N-c_i}{c_i(N+1)}} \leq\frac1{2\sqrt{c_i}}. \tag{3} \] If $c_i=0$, the total variation is at most one.
\end{lemma}

\begin{proof}
Work on the face indexed by $\supp(c)$. The ratio of Dirichlet densities is \[ \frac{d\D_{c+e_i}}{d\D_c}(p)=\frac{Np_i}{c_i}. \] If $P\sim\D_c$, then $P_i\sim\operatorname{Beta}(c_i,N-c_i)$ and \[ \operatorname{Var}\!\left(\frac{NP_i}{c_i}\right)=\frac{N-c_i}{c_i(N+1)}. \] The first inequality follows from the density ratio and Cauchy--Schwarz. The second is immediate. If $c_i=0$, the two laws can lie on different faces, so we use the trivial bound. When $c_i=N$, both laws are the same point mass and the displayed bound is zero.
\end{proof}

If a label appears $n_i$ times, applying Lemma~\ref{lem:tv} successively gives a cumulative change proportional to \[ \sum_{j=1}^{n_i}j^{-1/2}=O(\sqrt{n_i}). \] This controls stability. To control the comparator term with the same update, we next exploit a property specific to exponential weights.

\section{Exponential Size Bias and Be the Leader}

Fix the entire outcome sequence only for this argument. Draw one global collection $W_1,\ldots,W_T\stackrel{\mathrm{iid}}\sim\operatorname{Exp}(1)$ and define the weighted leader after observing round $t$ by \[ A_{t+1}=\frac{\sum_{s\leq t}W_se_{y_s}}{\sum_{s\leq t}W_s}. \tag{4} \] For every realization of the weights and every proper loss, $A_{t+1}$ minimizes the weighted cumulative loss through $t$. The standard be-the-leader lemma \citep{cesabianchi2006} therefore gives \[ \sum_{t=1}^TW_t\ell(A_{t+1},y_t) \leq\min_{p\in\Delta_K}\sum_{t=1}^TW_t\ell(p,y_t). \tag{5} \] Although the global weights couple rounds, we use them only to prove an inequality about the marginal Dirichlet laws in (BB). The actual algorithm resamples independently.

\begin{lemma}[Dirichlet be-the-perturbed-leader]
\label{lem:btpl}
For every fixed outcome sequence and every bounded proper loss, \[ \sum_{t=1}^T\E_{P\sim\D_{c_t}}\ell(P,y_t) -\min_{p\in\Delta_K}\sum_{t=1}^T\ell(p,y_t) \leq\sum_{i=1}^K\sum_{j=1}^{n_i}\frac1{\sqrt j}. \tag{6} \]
\end{lemma}

\begin{proof}
For $W\sim\operatorname{Exp}(1)$ and every integrable $f$, \[ \E[Wf(W)]=\E[f(W+W')], \tag{7} \] where $W'$ is an independent exponential variable. Indeed, the size-biased density $we^{-w}$ is $\operatorname{Gamma}(2,1)$, the law of $W+W'$.

Suppose $y_t=i$ and it is the $j$th occurrence of class $i$. Ordinarily, (4) has law $\D_{c_t}$. In the $W_t$-weighted expectation on the left of (5), identity (7) adds one exponential copy of the current observation. Thus its law is $\D_{c_t+e_i}$. Taking expectations in (5) gives
\begin{align*}
\sum_t\E_{\D_{c_t+e_{y_t}}}\ell(P,y_t)
&\leq\E\!\left[\min_p\sum_tW_t\ell(p,y_t)\right]\\
&\leq\min_p\sum_t\ell(p,y_t),
\end{align*}
where the last step uses $\E W_t=1$ and expectation of a minimum is at most the minimum of expectations. Since $\ell\in[-1,1]$, replacing $\D_{c_t+e_i}$ by $\D_{c_t}$ costs at most twice their total variation. Lemma~\ref{lem:tv} bounds that cost by $1/\sqrt j$. Summing proves (6).
\end{proof}

This is the main reason for exponential rather than generic random weights. The size-biased law is exactly one additional weight from the same family, and gamma additivity translates that observation-level identity into a one-count Dirichlet update.

\section{Proof of the Main Theorems}

For $c\neq0$, define \[ L_t(c)=\E_{P\sim\D_c}\ell(P,y_t). \] Let $\D_0$ be the point mass at $P_1$, so this notation also covers the first round. Add and subtract $L_t(c_t)$ to decompose regret as
\begin{align}
\Reg_\ell(T)
={}&\sum_{t=1}^T\bigl(L_t(c_{t-1})-L_t(c_t)\bigr) \tag{8}\\
&+\sum_{t=1}^TL_t(c_t)-\min_p\sum_{t=1}^T\ell(p,y_t).\nonumber
\end{align}
The second line is bounded by Lemma~\ref{lem:btpl}. For the first line, consider the $j$th occurrence of class $i$. If $j=1$, total variation is at most one. If $j\geq2$, Lemma~\ref{lem:tv} with the previous count $j-1$ gives \[ L_t(c_{t-1})-L_t(c_t)\leq\frac1{\sqrt{j-1}}. \] Consequently, \[ \Reg_\ell(T)\leq 2m+\sum_i\sum_{j=1}^{n_i-1}\frac1{\sqrt j} +\sum_i\sum_{j=1}^{n_i}\frac1{\sqrt j}. \] Using $\sum_{j=1}^n j^{-1/2}\leq2\sqrt n$ and $m\leq\sum_i\sqrt{n_i}$ proves \[ \Reg_\ell(T)\leq6\sum_i\sqrt{n_i}. \] Cauchy--Schwarz gives $\sum_i\sqrt{n_i}\leq\sqrt{mT}$, and $m\leq\min\{K,T\}$, proving Theorem~\ref{thm:worst}.

For Theorem~\ref{thm:smooth}, set $q=c_{t-1}/(t-1)$ for $t\geq2$. If $P\sim\D_{c_{t-1}}$, standard Dirichlet moments give \[ \E[P]=q, \qquad \E\|P-q\|_2^2=\frac{1-\|q\|_2^2}{t}\leq\frac1t. \tag{9} \] Smoothness and the zero mean of $P-q$ imply \[ \E[\ell(P,y_t)-\ell(q,y_t)] \leq\frac\beta2\E\|P-q\|_2^2 \leq\frac\beta{2t}. \] At round one, the bootstrap and FTL use the same initialization. Summing the comparison and adding the common comparator proves the transfer inequality. The prior $O(\beta\log T)$ FTL bound gives the corollary.

To prove Corollary~\ref{cor:finite}, condition at round $t$ on the history and the outcome selected before the fresh draw. For each fixed $\ell$, the difference between $\ell(P_t,y_t)$ and its conditional expectation is a martingale difference whose range has length two. Hoeffding's lemma and a union bound over $\mathcal L_0$ give (1). The deterministic count-vector inequality proved above bounds the sum of conditional expectations minus the realized-path comparator. The standard log-moment-generating-function argument gives (2). The supplement gives the details and an extension through uniform covers.

\section{Discussion}

The Bayesian bootstrap resolves the geometric obstacle in simultaneous U-calibration without constructing a new barrier or projection. Its randomness is automatically centered, remains on the current simplex face, and decays at the statistical $t^{-1/2}$ scale. The exponential representation supplies a separate online-learning advantage through exact size bias.

The count-adaptive bound may be useful beyond worst-case dimension. It depends on the realized occupancy profile rather than the nominal alphabet size. The proof also isolates a reusable design principle: seek random-weight families whose size-biased law equals the original law plus one independent weight, then combine weighted be-the-leader with stability of the induced empirical distribution.

There are two clear boundaries. First, our guarantee is per-loss expected regret rather than expected supremum over losses. Second, Dirichlet draws can approach the boundary, so relative smoothness with respect to the log barrier does not follow from the Euclidean variance argument. The self-concordant construction of \citet{frongillo2026} remains stronger for that nonsmooth class. Determining whether one proper forecaster can be rate-optimal for every individual proper loss remains open.

\paragraph{Reproducibility.} The supplement gives complete proofs, including face-degenerate cases and the adaptive protocol. A deterministic numerical audit integrates the Dirichlet update distance, checks exact moments, and solves the binary worst-sequence dynamic program for squared and V-shaped losses. No numerical calculation is used as a proof.

\bibliography{references}

\end{document}
